\documentclass[../../main.tex]{subfiles}
\begin{document}

\begin{flushright}
\footnotesize\textit{【自動文字起こし・要確認】}
\end{flushright}

\noindent
{\bf [解]}

単位行列を $I$ とおく．ケーリー・ハミルトンの定理から，
\begin{align}
A^2 - \Bigl(\dfrac{1}{3} + 3\Bigr) A + \det(A) \cdot I &= 0 \\
\implies A^2 &= \dfrac{10}{3}A - I \quad \cdots \text{①}
\end{align}
である．したがって，$n \in \mathbb{N}_{\ge 2}$ に対し，
\begin{align}
A^{n+1} &= A^{n-1} \cdot A^2 = A^{n-1} \Bigl( \dfrac{10}{3} A - I \Bigr) \\
&= \dfrac{10}{3} A^n - A^{n-1} \quad \cdots \text{②}
\end{align}
となる．ここで，数列 $c_n$ が $c_{n+2} = \dfrac{10}{3} c_{n+1} - c_n$ を満たすとき，その一般項は $p, q$ を定数として
\begin{align}
c_n = 3^n p + \Bigl(\dfrac{1}{3}\Bigr)^n q
\end{align}
で与えられることに注意する．

\medskip
\noindent
\paragraph{(1)}

$A = \begin{pmatrix} 1/3 & 5 \\ 0 & 3 \end{pmatrix}$，$A^2 = \begin{pmatrix} 1/9 & 50/9 \\ 0 & 9 \end{pmatrix}$ である．ここで
\begin{align}
A^n &= 3^{n-1} \begin{pmatrix} 0 & 45/8 \\ 0 & 3 \end{pmatrix} + \Bigl(\dfrac{1}{3}\Bigr)^{n-1} \begin{pmatrix} 1/3 & -5/8 \\ 0 & 0 \end{pmatrix} \\
&= \begin{pmatrix} (1/3)^n & (3^n - (1/3)^n)\dfrac{15}{8} \\ 0 & 3^n \end{pmatrix} \quad \cdots \text{③}
\end{align}
であることを帰納的に示す．$n=1,2$ の時は成立する．そこで $n=k, k+1 \ (k \in \mathbb{N})$ での③の成立を仮定すると，②から
\begin{align}
A^{k+2} &= \dfrac{10}{3} \begin{pmatrix} (1/3)^{k+1} & (3^{k+1} - (1/3)^{k+1})\dfrac{15}{8} \\ 0 & 3^{k+1} \end{pmatrix} - \begin{pmatrix} (1/3)^k & (3^k - (1/3)^k)\dfrac{15}{8} \\ 0 & 3^k \end{pmatrix} \\
&= \begin{pmatrix} (1/3)^{k+2} & (3^{k+2} - (1/3)^{k+2})\dfrac{15}{8} \\ 0 & 3^{k+2} \end{pmatrix}
\end{align}
となり，$n=k+2$ でも成立する．以上より示された．

したがって，
\begin{align}
A^n = \begin{pmatrix} (1/3)^n & (3^n - (1/3)^n)\dfrac{15}{8} \\ 0 & 3^n \end{pmatrix}
\end{align}
である．さらに
\begin{align}
\begin{pmatrix} a_n \\ b_n \end{pmatrix} &= A^n \begin{pmatrix} a \\ b \end{pmatrix} \\
&= \begin{pmatrix} (1/3)^n a + \dfrac{15}{8}(3^n - (1/3)^n) b \\ 3^n b \end{pmatrix} \quad \cdots \text{④}
\end{align}
だから，$\alpha = 3^n$ とおくと，
\begin{align}
a_n^2 + b_n^2 = \Bigl\{ \dfrac{15}{8} b \alpha + \Bigl(a - \dfrac{15}{8}b\Bigr) \dfrac{1}{\alpha} \Bigr\}^2 + (b\alpha)^2
\end{align}
であり，$n \to \infty$ ($\alpha \to \infty$) において，

\paragraph{$b \ne 0$ のとき}
\begin{align}
u &= \lim_{n \to \infty} \dfrac{a_n}{\sqrt{a_n^2 + b_n^2}} \\
&= \dfrac{\dfrac{15}{8}b}{\sqrt{\Bigl(\dfrac{15}{8}b\Bigr)^2 + b^2}} \\
&= \dfrac{15}{17}\dfrac{b}{|b|}
\end{align}
\begin{align}
v &= \lim_{n \to \infty} \dfrac{b_n}{\sqrt{a_n^2 + b_n^2}} \\
&= \dfrac{b}{\sqrt{\Bigl(\dfrac{15}{8}b\Bigr)^2 + b^2}} \\
&= \dfrac{8}{17}\dfrac{b}{|b|}
\end{align}

\paragraph{$b = 0$ のとき}
④に代入して $\begin{pmatrix} a_n \\ b_n \end{pmatrix} = \begin{pmatrix} a/\alpha \\ 0 \end{pmatrix}$ だから，$v=0, u = \dfrac{a}{|a|}$ となる．

\medskip
以上をまとめて，
\begin{align}
\begin{cases}
b=0 \text{ のとき } & (u,v) = \Bigl(\dfrac{a}{|a|}, 0\Bigr) \\[1ex]
b \ne 0 \text{ のとき } & (u,v) = \Bigl(\dfrac{15}{17}\dfrac{b}{|b|}, \dfrac{8}{17}\dfrac{b}{|b|}\Bigr)
\end{cases}
\end{align}
となる．

\newpage

\end{document}
